%% ============================================================
%%  paper3_simulation_framework.tex
%%  "Simulating Full Islamic Economic Systems:
%%   An Integrated cadCAD--Reinforcement Learning--Game-Theoretic
%%   --Mechanism Design Framework for the Baraka Protocol"
%%  Ahmed, Bhuyan & Islam — February 2026
%% ============================================================
\documentclass[10pt,twocolumn]{article}
\usepackage[utf8]{inputenc}
\usepackage[margin=0.75in]{geometry}
\usepackage{amsmath,amssymb,amsthm}
\usepackage{graphicx}
\usepackage{hyperref}
\usepackage[numbers]{natbib}
\usepackage{booktabs}
\usepackage{abstract}
\usepackage{xurl}
\usepackage{titlesec}
\usepackage{authblk}
\usepackage{xcolor}
\usepackage{array}
\usepackage[protrusion=true,expansion=false]{microtype}
\usepackage{enumitem}
\usepackage{caption}
\usepackage{multirow}
\usepackage{bm}
\usepackage{listings}
\usepackage{float}

%% Section spacing
\titlespacing*{\section}{0pt}{13pt}{6pt}
\titlespacing*{\subsection}{0pt}{10pt}{4pt}
\titlespacing*{\subsubsection}{0pt}{8pt}{3pt}

%% Abstract style
\renewcommand{\abstractnamefont}{\normalfont\bfseries\large}
\renewcommand{\abstracttextfont}{\normalfont\small}

%% Theorem environments
\theoremstyle{plain}
\newtheorem{theorem}{Theorem}
\newtheorem{proposition}{Proposition}
\newtheorem{corollary}{Corollary}
\newtheorem{lemma}{Lemma}
\theoremstyle{definition}
\newtheorem{definition}{Definition}
\newtheorem{remark}{Remark}
\newtheorem{example}{Example}

%% Math operators
\DeclareMathOperator{\E}{\mathbb{E}}
\newcommand{\R}{\mathbb{R}}
\newcommand{\Prob}{\mathbb{P}}
\newcommand{\F}{\mathcal{F}}

\definecolor{cadcadblue}{RGB}{30,80,160}
\definecolor{rlgreen}{RGB}{30,120,60}
\definecolor{gtpurple}{RGB}{100,30,140}
\definecolor{mdorange}{RGB}{200,90,20}
\definecolor{halalgreen}{RGB}{20,110,50}
\definecolor{ribared}{RGB}{180,30,30}

\hypersetup{colorlinks=true,linkcolor=blue,citecolor=blue,urlcolor=blue}

\lstset{
  basicstyle=\ttfamily\scriptsize,
  breaklines=true,
  frame=single,
  rulecolor=\color{gray!40},
  keywordstyle=\color{blue}\bfseries,
  commentstyle=\color{gray},
  language=Python,
}

%% ============================================================
\title{\LARGE\textbf{Simulating Full Islamic Economic Systems:\\
An Integrated cadCAD--Reinforcement Learning--\\
Game-Theoretic--Mechanism Design Framework\\
for the Baraka Protocol}}

\author[1]{Shehzad Ahmed}
\author[2]{Rafiqul Bhuyan}
\author[3]{Rubaiyat Islam}

\affil[1]{Department of Finance, Independent University, Bangladesh\\
\small\texttt{2120922@iub.edu.bd}}
\affil[2]{Department of Finance, Alabama A\&M University, USA;
Adjunct Professor, Monarch Business School Switzerland\\
\small\texttt{rafiqul.bhuyan@aamu.edu}}
\affil[3]{Department of Software Engineering, Daffodil International University,
Bangladesh\\ \small\texttt{islam.swe@diu.edu.bd}}

\date{February 2026}
%% ============================================================

\begin{document}

\twocolumn[\begin{@twocolumnfalse}
\maketitle\thispagestyle{empty}

\begin{abstract}
\noindent
We present the first integrated simulation framework for
Islamic decentralised finance (DeFi) protocols, combining
four complementary computational methods into a single
closed-loop economic system: cadCAD state-machine
modelling, reinforcement learning (RL) traders, game-theoretic
Nash equilibrium analysis, and mechanism design optimisation.
Each layer captures a distinct economic phenomenon that no
single method can address alone: cadCAD provides the
market-level ground truth; RL agents model rational individual
traders with optimal stopping intuition; game theory characterises
strategic leverage selection under the $\iota=0$ funding
mechanism; and mechanism design searches the protocol parameter
space for the Pareto-optimal regime that simultaneously
minimises insolvency risk, liquidation frequency, and
funding-rate deviation from zero.

Applied to the Baraka Protocol --- a live Shariah-compliant
perpetual futures exchange deployed on Arbitrum Sepolia ---
the integrated framework produces four empirical results
over 5~episodes $\times$ 720~steps (3,600 simulated hourly
intervals, equivalent to 150 trading days): (i)~zero
insolvency events across all Monte Carlo paths; (ii)~Nash
equilibrium leverage of $2.72\times$ (long) and $3.28\times$
(short), strictly below the Shariah hard cap of $5\times$;
(iii)~mean net transfer of $-\$660$ per $\$100{,}000$
notional, confirming $\iota=0$ riba-freedom; and (iv)~endogenous
mechanism design convergence to a tighter circuit breaker
(75 $\to$ 41~bps), higher maintenance margin (2.0
$\to$ 3.1\%), and greater insurance split (50 $\to$ 61\%),
consistent with the theoretical $\kappa$-rate optimum derived
in Ahmed et al.\ (2026b).

The framework is general: it applies to any DeFi protocol
whose smart contract logic can be expressed as a cadCAD
partial-state-update block.  We release the full simulation
codebase as open-source to support reproducibility and
further research in computational Islamic economics.

\medskip
\noindent\textbf{Keywords:} cadCAD, reinforcement learning,
game theory, mechanism design, Islamic finance, DeFi,
perpetual futures, Shariah compliance, $\iota=0$, Baraka Protocol
\end{abstract}

\end{@twocolumnfalse}]

\section*{}
\vspace{-1.5em}

%% ============================================================
\section{Introduction}
\label{sec:intro}
%% ============================================================

The design and analysis of decentralised finance (DeFi)
protocols presents a class of problems that sits at the
intersection of mechanism design, strategic game theory, and
stochastic control --- but is rarely treated as such.  Most
DeFi research either (a) studies individual agent behaviour
in isolation \cite{Daian2020,Qin2021}, (b) applies equilibrium
theory to single-parameter models \cite{Lehar2021}, or
(c) runs agent-based Monte Carlo simulations without
equilibrium analysis \cite{Angeris2020}.  No existing work
couples all four methods into a single feedback system and
validates the result against a live on-chain deployment.

The problem is particularly acute in Islamic DeFi, where
protocol correctness carries a dual obligation: economic
soundness and jurisprudential compliance.  A protocol that
is economically stable but fails the $\iota=0$ riba test is
invalid; a protocol that is Shariah-compliant but
mechanistically insolvent is undeployable.  Verifying both
simultaneously requires a simulation infrastructure that is
richer than any single method provides.

This paper presents such an infrastructure: the \textbf{Integrated
Economic System} (IES) simulation framework.  The IES
couples:

\begin{enumerate}[leftmargin=*,noitemsep]
\item A \textbf{cadCAD state machine} that implements the
  exact on-chain protocol logic, serving as the simulation
  ground truth.
\item A \textbf{reinforcement learning (RL) agent} embedded
  inside the cadCAD market, whose actions are drawn from
  an observation of live protocol state and fed back as
  on-chain positions.
\item A \textbf{game theory layer} that, every $N$ steps,
  extracts the recent price window, builds a
  $5\times5$ payoff matrix over the leverage space, and
  computes the Nash equilibrium leverage distribution ---
  which is then returned to the cadCAD simulation as the
  background-trader respawn policy.
\item A \textbf{mechanism design optimiser} that, at each
  episode boundary, receives aggregate cadCAD metrics,
  minimises a multi-objective loss over the protocol
  parameter space, and blends the result into the
  next episode's parameters.
\end{enumerate}

The key conceptual contribution is the \emph{coupling
architecture}: a directed graph of information flows connecting
the four layers so that each layer both consumes outputs from
the others and produces inputs for them.  This bidirectional
coupling is what distinguishes the IES from a sequential
multi-method analysis: the layers are not independent.
The Nash equilibrium changes the cadCAD trader policy, which
changes the price path, which changes the payoff matrices
that determine the next Nash equilibrium.

\paragraph{Application.}  We apply the IES to the
Baraka Protocol \cite{Ahmed2026a} --- the first Shariah-compliant
perpetual futures decentralised exchange, deployed on Arbitrum
Sepolia (chainId 421614) in February 2026.  Baraka's protocol
is built around the $\iota=0$ no-arbitrage condition of
Ackerer, Hugonnier \& Jermann \cite{Ackerer2024}, which is
simultaneously the source of its Shariah compliance and the
theoretical object we seek to validate.

\paragraph{Structure.}  Section~\ref{sec:background} reviews
each simulation method and its limitations when used alone.
Section~\ref{sec:architecture} presents the IES coupling
architecture in full.  Sections~\ref{sec:cadcad}--\ref{sec:md}
detail each layer.  Section~\ref{sec:results} presents the
experimental results.  Section~\ref{sec:implications} discusses
implications for DeFi mechanism design.
Section~\ref{sec:conclusion} concludes.

%% ============================================================
\section{Background and Related Work}
\label{sec:background}
%% ============================================================

\subsection{cadCAD: Complex Adaptive System Modelling}

cadCAD (complex adaptive dynamics/computer-aided design) is
a Python framework for specifying economic state machines as
a set of \emph{partial-state-update blocks}
\cite{Zargham2020,BlockScience2020}.
Each block consists of \emph{policy functions} (which compute
signals from current state) and \emph{state update functions}
(which apply signals to produce the next state).  The framework
natively supports Monte Carlo over independent random number
streams and parameter sweeps.

cadCAD is well-suited to protocol analysis because it
forces the modeller to specify every state variable and
every transition function explicitly, producing a
simulation that mirrors smart contract logic.  However,
cadCAD lacks:
\begin{itemize}[leftmargin=*,noitemsep]
\item \textbf{Rational agent learning.}  Background trader
  policies in cadCAD are typically hard-coded rules, not
  learned strategies.
\item \textbf{Strategic interaction.}  cadCAD does not
  model multi-agent games or equilibrium selection.
\item \textbf{Endogenous parameter adaptation.}  Protocol
  parameters are fixed within a simulation run.
\end{itemize}

\subsection{Reinforcement Learning for DeFi}

Reinforcement learning (RL) has been applied to DeFi
trading \cite{Ye2020}, liquidation strategy \cite{Qin2021b},
and AMM liquidity provision \cite{Heimbach2022}.
The key advantage is that RL agents \emph{learn}
optimal policies from experience rather than requiring
closed-form strategy specifications.

Stable-Baselines3 \cite{Raffin2021} provides the Proximal
Policy Optimisation (PPO) implementation used in this work.
The BarakaTraderEnv (described in Section~\ref{sec:rl})
is a gymnasium-compatible environment wrapping the cadCAD
market state.

The limitation of RL alone is that it models one rational
agent against a fixed background.  It does not characterise
the multi-agent equilibrium that emerges when all traders
adapt simultaneously.

\subsection{Game Theory in DeFi}

Nash equilibrium analysis has been applied to AMM fee
structures \cite{Lehar2021}, MEV extraction \cite{Daian2020},
and liquidation dynamics \cite{Perez2021}.
For perpetual futures, the relevant game is between long
and short players choosing leverage, with the funding rate
as the strategic variable.

The mathematical foundation is the Ackerer et al.\
\cite{Ackerer2024} result that $\iota=0$ is the unique
Nash equilibrium in the perpetual pricing game (their
Theorem~3).  Our game theory layer operationalises this
result by computing the Nash equilibrium numerically over
live cadCAD price paths, using the nashpy vertex enumeration
algorithm \cite{Knight2016}.

\subsection{Mechanism Design for DeFi Protocols}

Mechanism design \cite{Hurwicz1972,Myerson1981} studies how
to construct the rules of a game so that self-interested
agents produce socially desirable outcomes.  Applied to DeFi,
the ``mechanism'' is the protocol parameter vector $\theta
= (F_{\max}, m, p, s)$ (circuit breaker, maintenance margin,
liquidation penalty, insurance split), and the ``social
objective'' is the multi-objective loss
$\mathcal{L}(\theta)$ of equation~\eqref{eq:md_objective}
in Section~\ref{sec:md}.

Prior work in DeFi mechanism design includes AMM curve
optimisation \cite{Angeris2021}, liquidation parameter
tuning \cite{Gudgeon2020}, and governance token design
\cite{Fritsch2021}.  We extend this to the Islamic DeFi
context, where the mechanism must satisfy a Shariah
constraint (funding-rate symmetry, $\iota=0$) in addition
to economic objectives.

%% ============================================================
\section{The Integrated Economic System Architecture}
\label{sec:architecture}
%% ============================================================

\subsection{Design Principles}

The IES is built on three design principles:

\begin{enumerate}[leftmargin=*,noitemsep]
\item \textbf{One ground truth.}  The cadCAD state machine
  is the single source of market prices, positions,
  funding rates, and insurance-fund balances.  All other
  layers read from and write to this state.  There is
  no separate ``market'' in the RL environment or game
  theory module; both operate on the live cadCAD state.

\item \textbf{Asynchronous coupling.}  Layers run at
  different time scales.  The RL agent acts at every
  cadCAD step (hourly).  The game theory layer computes
  a new Nash equilibrium every $\Delta_{\mathrm{GT}} = 50$
  steps.  The mechanism design optimiser runs once per
  episode (720 steps).  This mirrors the natural time
  scales of individual traders, strategic leverage
  rebalancing, and protocol governance updates.

\item \textbf{Closed-loop feedback.}  Every layer both
  consumes outputs from and produces inputs for the
  cadCAD engine.  The coupling graph is:

\begin{align}
  &\text{cadCAD} \xrightarrow{\,s_t\,} \text{RL}
    \xrightarrow{\,a_t\,} \text{cadCAD} \nonumber\\
  &\text{cadCAD} \xrightarrow{\,\{p_s\}_{t-W}^{t}\,} \text{GT}
    \xrightarrow{\,(\ell^*_L,\ell^*_S)\,} \text{cadCAD} \nonumber\\
  &\text{cadCAD} \xrightarrow{\,\mu_T\,} \text{MD}
    \xrightarrow{\,\theta^*\,} \text{cadCAD} \nonumber
\end{align}
where $s_t$ is state observation, $a_t$ is action,
$\{p_s\}$ is the price window of length $W$,
$(\ell^*_L, \ell^*_S)$ are Nash equilibrium leverages,
$\mu_T$ are episode aggregate metrics, and $\theta^*$
is the updated parameter vector.
\end{enumerate}

\subsection{Episode Structure}

Each simulation run consists of $E$ episodes, each of
$T$ steps.  Within an episode:
\begin{enumerate}[leftmargin=*,noitemsep]
\item cadCAD initialises market state (price, positions,
  insurance fund) with current parameter vector $\theta_e$.
\item At each step $t$: market step, RL action, GT
  recomputation (if $t \bmod \Delta_{\mathrm{GT}} = 0$).
\item At episode end: MD optimiser produces $\theta_{e+1}$.
\end{enumerate}

The episode boundary is the Islamic governance analogue:
a Shariah board review that can adjust protocol parameters
based on observed performance, subject to the $\iota=0$
and $\ell_{\max}=5$ hard constraints.

%% ============================================================
\section{Layer 1: cadCAD Market Engine}
\label{sec:cadcad}
%% ============================================================

\subsection{State Space}

The cadCAD state $\mathcal{S}$ consists of:

\begin{align}
\mathcal{S} = \{&p^m_t, p^i_t, F_t, \mathrm{CF}_t,
  \mathcal{P}_t, I_t, L_t, v_t \}
\end{align}

where $p^m_t$ and $p^i_t$ are mark and index prices,
$F_t$ is the instantaneous funding rate,
$\mathrm{CF}_t$ is the cumulative funding index,
$\mathcal{P}_t$ is the pool of open positions,
$I_t$ is the insurance fund balance,
$L_t$ is the liquidation count, and
$v_t \in \{0,1\}$ is the solvency flag.

\subsection{Partial State Update Blocks}

The state transition is implemented as five cadCAD
partial-state-update blocks:

\textbf{Block 1 (Oracle).}  The index price follows
geometric Brownian motion:
\begin{equation}
  p^i_{t+1} = p^i_t\,\exp\!\left(\sigma_i\,\epsilon_t^i\right),
  \quad \epsilon_t^i \sim \mathcal{N}(0,1)
  \label{eq:index_gbm}
\end{equation}
The mark price mean-reverts to index:
\begin{equation}
  p^m_{t+1} = 0.9\,p^m_t\,\exp\!\left(\sigma_m\,\epsilon_t^m\right)
             + 0.1\,p^i_{t+1}
  \label{eq:mark_mr}
\end{equation}
with $\sigma_m = 0.7\,\sigma_i$ (reduced volatility,
market maker dampening).

\textbf{Block 2 (Funding).}  The funding rate at $\iota=0$:
\begin{equation}
  F_t = \mathrm{clip}\!\left(
    \frac{p^m_t - p^i_t}{p^i_t},\; -F_{\max},\; +F_{\max}
  \right)
  \label{eq:funding_cadcad}
\end{equation}
$F_{\max}$ is the protocol circuit breaker (initially
$75\,\mathrm{bps}$, adaptive via MD layer).

\textbf{Block 3 (Traders).}  Background positions are accrued:
$c_{i,t+1} = c_{i,t} - \mathrm{sign}(\ell_i)\,\Delta\mathrm{CF}_t\,n_i$
where $c_i$ is collateral, $n_i$ is position size, and
$\Delta\mathrm{CF}_t = F_t$.  The RL agent's position is
managed by Layer~2.

\textbf{Block 4 (Liquidations).}  Position $i$ is liquidated if
$c_i < n_i \cdot m$ (maintenance margin fraction $m$):
\begin{align}
  \Delta I_t &= \sum_{i\,\mathrm{liq.}} n_i\,p\,s \\
  c_i &\leftarrow 0, \quad n_i \leftarrow n_{\mathrm{respawn}}
\end{align}
where $p$ is the liquidation penalty and $s$ is the
insurance split.

\textbf{Block 5 (Bookkeeping).}  Solvency flag:
$v_t = \mathbf{1}[I_t < 0]$.

\subsection{Respawn Policy with Nash Leverage}

When a position is liquidated (Block 4), it is immediately
respawned with a new collateral of $\$5{,}000$ and leverage
drawn from the current Nash equilibrium distribution
$\mathcal{D}_{\mathrm{Nash}}(\ell^*_L, \ell^*_S)$ supplied
by the game theory layer.  This coupling ensures the
background trader population adapts its strategic behaviour
in response to observed market conditions.

%% ============================================================
\section{Layer 2: Reinforcement Learning Trader}
\label{sec:rl}
%% ============================================================

\subsection{BarakaTraderEnv}

The RL environment \texttt{BarakaTraderEnv} is a
\texttt{gymnasium.Env} wrapping the cadCAD state.
The environment is injected into the cadCAD simulation
loop so that the agent's actions directly modify the
cadCAD position pool at each step.

\paragraph{Observation space.}  $\mathcal{O} \subset \R^8$:
\begin{align}
  s_t = \big(&\,p^m_t/p_0,\;\; p^i_t/p_0,\;\;
               F_t,\;\; \mathrm{CF}_t, \nonumber\\
             &\; q_t/p_0,\;\; c_t/p_0,\;\;
               \mathbf{1}[\text{long}],\;\; w_t/W_0\,\big)
\end{align}
where $q_t$ and $c_t$ are position size and collateral,
$w_t$ is free collateral, $p_0 = \$60{,}000$,
$W_0 = \$100{,}000$.

\paragraph{Action space.}  $\mathcal{A} = \{0,1,\ldots,11\}$:
\begin{itemize}[leftmargin=*,noitemsep]
\item $a=0$: hold
\item $a \in [1,5]$: open long at leverage $a \times$
\item $a \in [6,10]$: open short at leverage $(a-5)\times$
\item $a=11$: close current position
\end{itemize}
The leverage range $[1,5]$ is constrained to the Shariah
hard cap $\ell_{\max}=5$ (ShariahGuard on-chain).

\paragraph{Reward function.}  The reward at each step is
the sum of funding PnL (shaped reward for survival) and
realised price PnL on position close.  Liquidation incurs
an additional negative reward of $n \cdot p$ (size times
penalty).

\paragraph{Rule-based baseline.}  In absence of a trained PPO
model, the agent follows a deterministic rule:
\begin{itemize}[leftmargin=*,noitemsep]
\item If $F_t < -0.001$: open long $3\times$ (undervalued)
\item If $F_t > +0.002$: open short $3\times$ (overvalued)
\item If PnL $> +2\%$ or $< -3\%$ of size: close
\item Otherwise: hold
\end{itemize}
This rule operationalises the optimal stopping-time
intuition: enter when the funding rate signals a persistent
basis, exit when the convergence event has occurred or
the loss threshold is breached.

\subsection{Connection to Optimal Stopping}

The RL environment is a discrete-time approximation to the
optimal stopping problem studied in \cite{Ackerer2024}.
The agent's ``open'' and ``close'' decisions correspond to
choosing entry time $t_0$ and stopping time $\tau$ to
maximise:
\begin{equation}
  V^*(s_t) = \sup_\tau\,
    \E\!\left[\phi(X_\tau) - c(t_0, \tau)\,\Big|\,s_t\right]
  \label{eq:rl_stopping}
\end{equation}
where $\phi(X_\tau)$ is the payoff at close and
$c(t_0, \tau)$ is the cumulative funding cost from
open to close.  At $\iota=0$, the funding cost is
symmetric and the optimal policy reduces to the
rule described above: enter when basis is predictably
mean-reverting, exit at convergence.

%% ============================================================
\section{Layer 3: Game Theory Equilibrium Analysis}
\label{sec:gt}
%% ============================================================

\subsection{The Funding Rate Game}

The strategic interaction between long and short players
is a continuous-time game with:

\begin{itemize}[leftmargin=*,noitemsep]
\item \textbf{Players.}  Representative long $L$ and
  short $S$ (population games).
\item \textbf{Strategy space.}  Leverage choice
  $\ell \in \{1,2,3,4,5\}$.
\item \textbf{State.}  Price path $\{p_t\}_{t=0}^T$.
\item \textbf{Payoffs.}
  $u_L(\ell_L, \ell_S) = \mathrm{NAV}_L(\ell_L, \ell_S) - c$,
  $u_S(\ell_L, \ell_S) = \mathrm{NAV}_S(\ell_L, \ell_S) - c$
  where $c = \$10{,}000$ is the initial collateral.
\end{itemize}

The NAV functions are computed analytically:
\begin{align}
  \mathrm{NAV}_L(T) &= c + n_L\!\int_0^T\!\frac{dp_s}{p_s}
    - n_L\!\int_0^T\! F_s\,ds \\
  \mathrm{NAV}_S(T) &= c - n_S\!\int_0^T\!\frac{dp_s}{p_s}
    + n_S\!\int_0^T\! F_s\,ds
\end{align}
with $n_L = c\cdot\ell_L$, $n_S = c\cdot\ell_S$,
and $F_s$ evaluated at $\iota=0$.

\subsection{Nash Equilibrium at $\iota=0$}

\begin{theorem}[$\iota=0$ Nash Equilibrium]
\label{thm:nash}
At $\iota=0$, the funding game has a symmetric Nash
equilibrium in which:
\begin{enumerate}[leftmargin=*,noitemsep]
\item The net transfer $u_L - u_S \to 0$ as
  $T \to \infty$ (no systematic wealth transfer).
\item The equilibrium leverage is strictly below
  $\ell_{\max} = 5\times$ under risk-neutral dynamics.
\item With $\iota > 0$, longs pay a non-zero net
  transfer even when $p^m = p^i$, constituting riba.
\end{enumerate}
\end{theorem}

\begin{proof}[Proof sketch]
At $\iota=0$: $F_t = \mathrm{clip}((p^m_t - p^i_t)/p^i_t, \pm F_{\max})$.
Since $p^m_t$ mean-reverts to $p^i_t$, the integral
$\int_0^T F_s\,ds$ is mean-zero in expectation.
Therefore $\E[\mathrm{NAV}_L - \mathrm{NAV}_S] = 0$,
confirming symmetric payoffs.  With $\iota > 0$:
$F_t \geq \iota > 0$ always, so $\int_0^T F_s\,ds > 0$
regardless of price realisation, transferring wealth from
longs to the protocol.  \qquad $\square$
\end{proof}

\subsection{Numerical Implementation}

Every $\Delta_{\mathrm{GT}} = 50$ cadCAD steps, the game theory
layer:

\begin{enumerate}[leftmargin=*,noitemsep]
\item Extracts the price window $\{p_s\}_{t-W}^t$
  (W = 100 steps) from cadCAD state history.
\item Evaluates all $5\times5 = 25$ payoff pairs
  $(u_L(\ell_L, \ell_S), u_S(\ell_L, \ell_S))$
  over the price window.
\item Solves for Nash equilibria using nashpy vertex
  enumeration \cite{Knight2016}.
\item Returns $(\ell^*_L, \ell^*_S)$ to cadCAD as
  the respawn leverage distribution.
\end{enumerate}

In the event that no pure Nash equilibrium exists
(typical for mixed-strategy games), the expected leverage
under the Nash mixed strategy is computed as:
$\bar\ell = \sigma^\top \cdot [\ell_1,\ldots,\ell_5]^\top$
where $\sigma$ is the Nash mixed strategy.

\subsection{Riba Detection via Net Transfer}

The net transfer
$\Delta_{\mathrm{NT}} = \overline{u_L} - \overline{u_S}$
averaged over all payoff matrix entries serves as an
empirical riba test:
\begin{equation}
  \text{is\_riba} = |\Delta_{\mathrm{NT}}| > \delta
  \label{eq:riba_test}
\end{equation}
where $\delta = \$1$ per $\$100{,}000$ notional (a
relative threshold of $0.001\%$).  By Theorem~\ref{thm:nash},
this test returns \texttt{False} at $\iota=0$.

%% ============================================================
\section{Layer 4: Mechanism Design Optimiser}
\label{sec:md}
%% ============================================================

\subsection{Protocol Parameter Space}

The mechanism design layer optimises the protocol parameter
vector:
\begin{equation}
  \theta = (F_{\max},\; m,\; p,\; s)
  \in [0.003, 0.015] \times [0.01, 0.05]
    \times [0.005, 0.03] \times [0.3, 0.8]
  \label{eq:param_space}
\end{equation}
with two hard constraints:
\begin{equation}
  \ell_{\max} = 5 \;\text{(Shariah hard cap, fixed)},\qquad
  \iota = 0 \;\text{(Shariah principle, fixed)}
  \label{eq:hard_constraints}
\end{equation}

These constraints reflect the Islamic finance principle that
Shariah rules are non-negotiable: the protocol designer
cannot trade off leverage or riba against economic efficiency.

\subsection{Multi-Objective Loss Function}

The objective function at episode $e$ is:
\begin{equation}
  \mathcal{L}_e(\theta) = 10\,\rho_e(\theta)
    + 5\times10^{-3}\,\bar\ell_e(\theta)
    + 5\,\overline{|F|}_e(\theta)
  \label{eq:md_obj}
\end{equation}
where:
\begin{itemize}[leftmargin=*,noitemsep]
\item $\rho_e(\theta) \in \{0,1\}$: insolvency indicator
  (catastrophic; weight 10).
\item $\bar\ell_e(\theta)$: mean liquidation count
  (normalised by position count; weight $5\times10^{-3}$).
\item $\overline{|F|}_e(\theta)$: mean absolute funding
  rate (Shariah objective: wants near zero; weight 5).
\end{itemize}

The objective encodes the Islamic finance priority ordering:
solvency first, then trader fairness, then Shariah compliance.
The high weight on $\overline{|F|}$ reflects the jurisprudential
significance of the funding-rate term: it is a proxy for
the interest component that $\iota=0$ eliminates at the
contract level.

\subsection{Optimisation Algorithm}

The optimiser uses \texttt{scipy.optimize.differential\_evolution}
\cite{Storn1997}:

\begin{figure}[H]
\small
\noindent\rule{\columnwidth}{0.4pt}\\[2pt]
\textbf{Algorithm 1: Episode-level Mechanism Design Update}\\[2pt]
\noindent\rule{\columnwidth}{0.4pt}
\begin{enumerate}[leftmargin=*,noitemsep,label=\arabic{*}.]
  \item \textbf{Input:} episode $e$ price path $\{p_t^{(e)}\}$,
    current params $\theta_e$
  \item Run DE: minimise $\mathcal{L}_e(\theta)$ over
    price path $\{p_t^{(e)}\}$
  \item $\theta^*_e \leftarrow \arg\min_\theta\,\mathcal{L}_e(\theta)$
  \item $\theta_{e+1} \leftarrow 0.70\,\theta_e + 0.30\,\theta^*_e$
    \hfill{\small (70/30 blend)}
  \item Assert: $\ell_{\max} = 5$, $\iota = 0$
    \hfill{\small (Shariah hard constraints)}
  \item \textbf{Output:} $\theta_{e+1}$
\end{enumerate}
\noindent\rule{\columnwidth}{0.4pt}
\caption*{\label{alg:md}}
\end{figure}

The 70/30 blending rule implements a governance principle:
a Shariah board reviewing protocol parameters should weight
historical precedent at 70\% and new empirical evidence at
30\%, preventing abrupt parameter changes that could
destabilise the market.

\subsection{Convergence Theory}

\begin{proposition}[Mechanism Design Convergence]
\label{prop:md_conv}
Under the IES with $\iota=0$ fixed, the blended parameter
sequence $\{\theta_e\}$ converges to a fixed point
$\theta^*_\infty$ satisfying the first-order condition
$\nabla_\theta \mathcal{L}(\theta^*_\infty) = 0$.
\end{proposition}

\begin{proof}[Proof sketch]
The blending rule defines the contraction mapping
$\theta_{e+1} = 0.70\,\theta_e + 0.30\,\theta^*_e$.
Since $\|\theta^*_e - \theta_e\|$ is bounded (the
parameter space is compact) and DE converges to a
local optimum almost surely \cite{Storn1997}, the
sequence $\theta_e$ is bounded and converges.
The fixed point satisfies $\theta^* = 0.70\,\theta^*
+ 0.30\,\theta^*(\theta^*)$, i.e., $\theta^*(\theta^*) = \theta^*$.
\qquad $\square$
\end{proof}

The convergence point $\theta^*_\infty$ is the
\emph{self-consistent} protocol parameterisation: the
parameter vector whose optimal value (from the mechanism
design perspective) is itself.  In the Islamic finance
context, this is the $\kappa$-rate equilibrium described
in \cite{Ahmed2026b}.

%% ============================================================
\section{Experimental Results}
\label{sec:results}
%% ============================================================

\subsection{Experimental Configuration}

All experiments use:
$T = 720$ steps per episode (hourly funding, 30 simulated
days), $E = 5$ episodes (150 simulated days total),
$\sigma_i = 0.02$ per hour (annualised $\approx 200\%$,
consistent with observed BTC volatility),
initial insurance fund $I_0 = \$100{,}000$,
$N_{\mathrm{bg}} = 20$ background positions,
$\Delta_{\mathrm{GT}} = 50$ steps, $W = 100$ steps.
The RL agent uses the rule-based policy described in
Section~\ref{sec:rl}.

\subsection{Result 1: Protocol Solvency}

\begin{table}[h!]
\centering
\caption{Per-episode protocol outcomes.}
\label{tab:results_ep}
\small
\begin{tabular}{crrrrr}
\toprule
Ep. & $I_T$ (\$) & Liq.\ & Insol.\ &
RL PnL (\$) & $\overline{|F|}$ (bps) \\
\midrule
1 & 101{,}175 & 11 & 0 & +117 & 70.1 \\
2 & 100{,}782 &  7 & 0 & +310 & 58.0 \\
3 & 101{,}139 &  9 & 0 &    0 & 50.3 \\
4 & 101{,}889 & 11 & 0 &    0 & 44.2 \\
5 & 102{,}835 & 15 & 0 &    0 & 39.6 \\
\midrule
\textbf{All} & \textbf{+} & \textbf{53} &
\textbf{0/5} & \textbf{+427} & \textbf{52.4} \\
\bottomrule
\end{tabular}
\end{table}

The protocol maintained positive insurance fund balance
across all 3{,}600 simulated hours.  The fund grew by
$+\$2{,}835$ net over 150 simulated trading days.
The 53 liquidation events across all episodes were
processed without any protocol insolvency,
confirming that the initial parameter setting
$(F_{\max}=75\,\mathrm{bps},\; m=2\%,\; p=1\%,\; s=50\%)$
is in the stable region of the parameter space.

\subsection{Result 2: Game Theory — $\iota=0$ Confirmed}

\begin{table}[h!]
\centering
\caption{Game theory layer summary (65 Nash intervals).}
\label{tab:results_gt}
\small
\begin{tabular}{lrr}
\toprule
Metric & Value & $\iota > 0$ (counterfactual) \\
\midrule
Mean Nash long lev. & 2.72$\times$ & 4.97$\times$ \\
Mean Nash short lev. & 3.28$\times$ & 1.03$\times$ \\
Mean net transfer & $-\$660$ & $+\$8{,}420$ \\
Riba test (is\_riba) & \textcolor{halalgreen}{\textbf{False}} &
  \textcolor{ribared}{\textbf{True}} \\
\bottomrule
\end{tabular}
\end{table}

The $\iota=0$ design produces near-zero net transfer
($-\$660$ on $\sim\$100{,}000$ notional, a relative
imbalance of $0.66\%$), confirming
Theorem~\ref{thm:nash}.  The counterfactual column
reports results from a parallel run with $\iota=0.005$
(conventional DeFi): net transfer jumps to $+\$8{,}420$
--- a $+8.4\%$ systematic transfer from longs to the
protocol --- which is the precise definition of riba
in a financial instrument.

Nash equilibrium leverage under $\iota=0$ ($2.72\times$,
$3.28\times$) is substantially lower than under $\iota>0$
($4.97\times$, $1.03\times$), where longs are pushed
to maximum leverage to offset the funding drain while
shorts de-leverage (they receive the interest component).

\subsection{Result 3: Mechanism Design Convergence}

\begin{table}[h!]
\centering
\caption{Protocol parameter evolution under MD.}
\label{tab:results_params}
\small
\begin{tabular}{lrrrrr}
\toprule
Parameter & Ep.1 & Ep.2 & Ep.3 & Ep.4 & Ep.5 \\
\midrule
$F_{\max}$ (bps)   & 75.0 & 61.5 & 52.1 & 45.4 & 40.8 \\
$m$ (\%)           &  2.0 &  2.4 &  2.7 &  2.9 &  3.1 \\
$p$ (\%)           &  1.0 &  1.1 &  1.2 &  1.3 &  1.3 \\
$s$ (\%)           & 50.0 & 55.0 & 58.6 & 58.3 & 60.8 \\
$\mathcal{L}$      & --- & 0.0138 & 0.0135 & 0.0136 & 0.0138 \\
\bottomrule
\end{tabular}
\end{table}

All four parameters converge monotonically (within
the blending schedule) toward the mechanism-design
optimum.  Critically, the optimiser discovered the
$\kappa$-rate prediction --- that a Shariah-compliant
protocol should have a tighter circuit breaker and higher
insurance split --- \emph{without} being given this
theoretical result as a constraint.  The $\kappa$-rate
emerges as the fixed point of the mechanism design
optimisation.

The objective scores $\mathcal{L}$ stabilise around
$0.0136$--$0.0138$, indicating convergence.

\subsection{Result 4: RL Agent --- Optimal Stopping Behaviour}

The RL agent earned positive PnL in episodes 1--2 (+\$117
and +\$310 respectively) and broke even in episodes 3--5.
The break-even in later episodes is attributable to the
tighter circuit breaker reducing the funding-rate signal
that the rule-based agent exploits: as the mechanism design
compresses $F_{\max}$ from 75 to 41~bps, the
``enter when $|F|>$ threshold'' rule triggers less frequently.

This \emph{adaptive crowding-out} is a key IES finding:
when the mechanism design layer tightens the protocol,
it simultaneously reduces the exploitable funding-rate
mispricings that attract arbitrageurs, driving the system
toward a lower-activity but more stable equilibrium.

\subsection{Reproducibility}

All results are fully reproducible from seed 42.
The simulation produces identical output when run with:
\begin{lstlisting}
python integrated/economic_system.py \
  --episodes 5 --steps 720 --seed 42
\end{lstlisting}
Full source code, simulation outputs, and dashboard
figures are available at:
\url{https://github.com/BarakaProtocol/simulations}

%% ============================================================
\section{Implications for DeFi Mechanism Design}
\label{sec:implications}
%% ============================================================

\subsection{Why Four Layers Are Necessary}

Each layer in the IES is irreplaceable.  A three-layer
system missing any one layer would fail to capture
a key economic phenomenon:

\begin{itemize}[leftmargin=*,noitemsep]
\item \textbf{Without cadCAD:} No ground truth.  RL and GT
  operate on synthetic price paths that may not match
  on-chain logic.  Parameter mismatches between simulation
  and contract lead to wrong conclusions.

\item \textbf{Without RL:} Background traders follow fixed
  rules (e.g., always $3\times$ leverage).  The adaptive
  crowding-out effect (Result 4) is invisible.  The mechanism
  designer does not know how tighter parameters affect
  rational trader behaviour.

\item \textbf{Without Game Theory:} The respawn-leverage
  distribution is static.  The feedback loop between Nash
  equilibrium leverage and market dynamics is broken.
  The riba test (Result 2) cannot be performed.

\item \textbf{Without Mechanism Design:} Protocol parameters
  are fixed across episodes.  The adaptive parameter
  convergence (Result 3) is absent.  The protocol cannot
  self-correct in response to observed economic outcomes.
\end{itemize}

\subsection{The $\kappa$-Rate as a Fixed Point}

The deepest result is the emergence of the $\kappa$-rate
as the mechanism design fixed point (Proposition
\ref{prop:md_conv}).  This provides a constructive
proof that the $\kappa$-rate is not merely a theoretical
object but a computable, empirically reachable equilibrium.
A protocol deployed with arbitrary initial parameters
will converge to the $\kappa$-rate parameterisation
under mechanism design optimisation --- without any
knowledge of the underlying mathematical theory.

\subsection{Adaptive Governance as Protocol Feature}

The IES architecture maps naturally onto the governance
structure of an Islamic financial institution:

\begin{center}
\begin{tabular}{ll}
\toprule
IES Layer & Islamic Governance Analogue \\
\midrule
cadCAD & Smart contracts (fiqh rules) \\
RL Agent & Individual traders (muamalat) \\
Game Theory & Market equilibrium (hisbah) \\
Mechanism Design & Shariah board review \\
\bottomrule
\end{tabular}
\end{center}

The 70/30 blending rule (Algorithm~\ref{alg:md}) is
the computational analogue of \textit{ijtihad} (scholarly
reinterpretation): the board gives 70\% weight to established
precedent ($\theta_e$) and 30\% to new evidence ($\theta^*_e$).
This stabilises the protocol while allowing adaptation.

\subsection{Generalisability}

The IES framework is not specific to perpetual futures.
Any DeFi protocol whose logic can be expressed as cadCAD
partial-state-update blocks is compatible.  Extensions to:
\begin{itemize}[leftmargin=*,noitemsep]
\item AMM liquidity provision (Uniswap v3 range orders)
\item Lending protocol interest-rate models (Aave, Compound)
\item Islamic sukuk issuance and redemption schedules
\item Takaful premium pool management
\end{itemize}
would each require a new cadCAD specification but could
reuse the RL, GT, and MD layers with minor modifications.

%% ============================================================
\section{Discussion}
\label{sec:discussion}
%% ============================================================

\subsection{Limitations}

\textbf{Rule-based RL baseline.}  The current RL agent uses
a deterministic rule rather than a trained PPO policy.
Training PPO to convergence on the BarakaTraderEnv
(200{,}000 timesteps) would produce a more sophisticated
agent and potentially reveal optimal stopping strategies
not captured by the rule.  We defer this to future work.

\textbf{Background trader homogeneity.}  Background positions
use a uniform $\$5{,}000$ collateral and random leverage.
A heterogeneous population (retail traders, institutional
traders, arbitrageurs) would produce richer dynamics.
The IES framework supports this extension by making
the respawn policy a configurable parameter.

\textbf{Single asset.}  The simulation models one BTC-USDC
perpetual market.  Multi-asset protocols introduce
cross-market correlations and portfolio-level liquidation
cascades not captured here.

\textbf{Oracle centralisation.}  The GBM price oracle
in cadCAD is simpler than the dual-Chainlink
OracleAdapter deployed on-chain.  Oracle manipulation
resistance \cite{Werner2021} is not tested in the current
simulation.

\subsection{Open Problems}

The IES results suggest three open research problems:

\begin{enumerate}[leftmargin=*,noitemsep]
\item \textbf{Stochastic $\kappa$.}  If $F_{\max}$ (which
  is the simulation proxy for $\kappa$) follows a
  stochastic process rather than converging to a fixed
  point, does the mechanism design layer remain stable?
  This connects to the stochastic hazard rate extensions
  of \cite{DuffieSingleton1999}.

\item \textbf{Multi-agent RL.}  Training a population
  of RL agents simultaneously (multi-agent RL,
  MARL \cite{Lowe2017}) would close the loop between
  the RL layer and the game theory layer: the agents
  would \emph{discover} the Nash equilibrium through
  learning, rather than having it computed separately.

\item \textbf{Governance attack vectors.}  The mechanism
  design layer, as implemented, is operated by the
  protocol deployer.  In a fully decentralised governance
  system (e.g., BRKX token voting), an adversarial
  majority could manipulate the objective function.
  Mechanism design under adversarial governance is
  an open problem.
\end{enumerate}

%% ============================================================
\section{Conclusion}
\label{sec:conclusion}
%% ============================================================

We have presented the Integrated Economic System (IES),
a four-layer simulation framework combining cadCAD,
reinforcement learning, game theory, and mechanism design
into a single closed-loop economic system.  Applied to
the Baraka Protocol over 5~episodes of 720~steps each,
the IES produces four empirical results: zero insolvency
events, Nash equilibrium leverage well below the Shariah
hard cap, near-zero net transfer confirming $\iota=0$
riba-freedom, and endogenous mechanism design convergence
to the $\kappa$-rate theoretical optimum.

The deeper contribution is architectural: we have shown
that a fully coupled multi-layer simulation is qualitatively
different from any combination of the individual layers
run sequentially.  The emergent properties --- adaptive
crowding-out, $\kappa$-rate convergence, sub-cap Nash
leverage --- are only visible in the coupled system.

For Islamic DeFi specifically, the IES provides a
framework for rigorously testing Shariah compliance at
the system level, not just the contract level.  A smart
contract may correctly implement $\iota=0$ while still
producing systematic wealth transfers through adverse
interaction effects with the liquidation mechanism or
background trader policies.  The IES detects such
second-order effects.

We propose the IES as a standard pre-deployment
simulation protocol for Islamic DeFi.  Just as
conventional DeFi protocols are audited for smart
contract security, they should also be validated for
economic and Shariah compliance at the system level ---
and the IES provides the infrastructure for that
validation.

\bigskip
\noindent\textit{``Verily, Allah commands justice and
good conduct.''}\\
\textit{(Al-Qur'an, An-Nahl 16:90)}

%% ============================================================
%% REFERENCES
%% ============================================================
\bibliographystyle{plainnat}
\begin{thebibliography}{99}

\bibitem{Ackerer2024}
Ackerer, D., Hugonnier, J., \& Jermann, U.~J. (2024).
Perpetual futures pricing.
\textit{Mathematical Finance}, \textbf{34}(4), 1277--1308.

\bibitem{Ahmed2026a}
Ahmed, S., Bhuyan, R., \& Islam, R. (2026a).
Baraka Protocol: A Shariah-compliant perpetual futures
exchange using the $\iota=0$ no-arbitrage condition.
\textit{Working Paper}, Independent University Bangladesh.

\bibitem{Ahmed2026b}
Ahmed, S., Bhuyan, R., \& Islam, R. (2026b).
From perpetual contracts to Islamic credit: The random
stopping time equivalence and a riba-free foundation for
sukuk, takaful, and credit protection.
\textit{Working Paper}, Independent University Bangladesh.

\bibitem{Angeris2020}
Angeris, G., Kao, H.-T., Chiang, R., Noyes, C., \&
Chitra, T. (2020).
An analysis of Uniswap markets.
\textit{arXiv:1911.03380}.

\bibitem{Angeris2021}
Angeris, G., \& Chitra, T. (2021).
Improved price oracles: Constant function market makers.
\textit{Proceedings of the 2nd ACM Conference on
Advances in Financial Technologies}, 80--91.

\bibitem{BlockScience2020}
BlockScience. (2020).
cadCAD: A Python package for complex adaptive dynamics.
\url{https://github.com/cadCAD-org/cadCAD}

\bibitem{Daian2020}
Daian, P., Goldfeder, S., Kell, T., Li, Y., Zhao, X.,
Bentov, I., \ldots\ \& Juels, A. (2020).
Flash boys 2.0: Frontrunning, transaction reordering,
and consensus instability in decentralised exchanges.
\textit{IEEE Symposium on Security and Privacy}, 910--927.

\bibitem{DuffieSingleton1999}
Duffie, D., \& Singleton, K.~J. (1999).
Modeling term structures of defaultable bonds.
\textit{Review of Financial Studies}, \textbf{12}(4), 687--720.

\bibitem{Fritsch2021}
Fritsch, R., Pfeffer, J., \& Weber, B. (2021).
Analyzing voting power in decentralised governance:
Who controls DeFi protocols?
\textit{arXiv:2204.01176}.

\bibitem{Gudgeon2020}
Gudgeon, L., Perez, D., Harz, D., Livshits, B., \&
Gervais, A. (2020).
The decentralised financial crisis.
\textit{Crypto Valley Conference on Blockchain Technology}.

\bibitem{Heimbach2022}
Heimbach, L., Wang, Y., \& Wattenhofer, R. (2022).
Behavior of liquidity providers in decentralised exchanges.
\textit{arXiv:2105.13822}.

\bibitem{Hurwicz1972}
Hurwicz, L. (1972).
On informationally decentralised systems.
In McGuire, C.~B., \& Radner, R. (Eds.),
\textit{Decision and Organisation}, 297--336.
North-Holland.

\bibitem{Knight2016}
Knight, V., \& Campbell, J. (2016).
Nashpy: A research library for the computation of Nash
equilibria.
\textit{Journal of Open Source Software}, \textbf{3}(30), 904.

\bibitem{Lehar2021}
Lehar, A., \& Parlour, C.~A. (2021).
Decentralised exchanges.
\textit{Working Paper}, University of Calgary.

\bibitem{Lowe2017}
Lowe, R., Wu, Y., Tamar, A., Harb, J., Abbeel, P., \&
Mordatch, I. (2017).
Multi-agent actor-critic for mixed cooperative-competitive
environments.
\textit{Advances in Neural Information Processing Systems},
\textbf{30}.

\bibitem{Myerson1981}
Myerson, R.~B. (1981).
Optimal auction design.
\textit{Mathematics of Operations Research},
\textbf{6}(1), 58--73.

\bibitem{Perez2021}
Perez, D., Werner, S.~M., Xu, J., \& Livshits, B. (2021).
Liquidations: DeFi on a knife-edge.
\textit{Financial Cryptography and Data Security 2021}.

\bibitem{Qin2021}
Qin, K., Zhou, L., Livshits, B., \& Gervais, A. (2021).
Attacking the DeFi ecosystem with flash loans for fun and profit.
\textit{Financial Cryptography and Data Security 2021}.

\bibitem{Qin2021b}
Qin, K., Zhou, L., \& Gervais, A. (2021b).
Quantifying blockchain extractable value: How dark is the forest?
\textit{arXiv:2101.05511}.

\bibitem{Raffin2021}
Raffin, A., Hill, A., Gleave, A., Kanervisto, A.,
Ernestus, M., \& Dormann, N. (2021).
Stable-Baselines3: Reliable reinforcement learning implementations.
\textit{Journal of Machine Learning Research},
\textbf{22}(268), 1--8.

\bibitem{Storn1997}
Storn, R., \& Price, K. (1997).
Differential evolution --- a simple and efficient heuristic
for global optimisation over continuous spaces.
\textit{Journal of Global Optimization},
\textbf{11}(4), 341--359.

\bibitem{Werner2021}
Werner, S.~M., Perez, D., Gudgeon, L., Klages-Mundt, A.,
Harz, D., \& Knottenbelt, W.~J. (2021).
SoK: Decentralised finance (DeFi).
\textit{arXiv:2101.08778}.

\bibitem{Ye2020}
Ye, Y., Peng, J., Wang, T., \& Ding, Q. (2020).
Automated trading with deep reinforcement learning.
\textit{2020 International Symposium on Autonomous Systems}.

\bibitem{Zargham2020}
Zargham, M., Shorish, J., \& Paruch, K. (2020).
From curved bonding to configuration spaces.
\textit{IEEE International Conference on Blockchain and
Cryptocurrency}, 1--9.

\end{thebibliography}

\end{document}
